\documentclass[10pt,letterpaper]{article}
\usepackage[letterpaper,margin=0.72in,headheight=16pt,footskip=28pt]{geometry}
\usepackage[T1]{fontenc}
\usepackage{lmodern}
\usepackage{microtype}
\usepackage[dvipsnames,table]{xcolor}
\usepackage{graphicx}
\usepackage{booktabs}
\usepackage{tabularx}
\usepackage{longtable}
\usepackage{array}
\usepackage{ragged2e}
\usepackage{pdflscape}
\usepackage{enumitem}
\usepackage{parskip}
\usepackage{needspace}
\usepackage[most]{tcolorbox}
\usepackage{listings}
\usepackage{upquote}
\usepackage{textcomp}
\usepackage{fancyhdr}
\usepackage{titlesec}
\usepackage{lastpage}
\usepackage{xurl}
\usepackage[hidelinks]{hyperref}
\usepackage{bookmark}

\definecolor{Navy}{HTML}{123047}
\definecolor{Teal}{HTML}{0E6670}
\definecolor{LightTeal}{HTML}{EAF4F4}
\definecolor{Critical}{HTML}{A12828}
\definecolor{High}{HTML}{A45D00}
\definecolor{Standard}{HTML}{356A45}
\definecolor{Defense}{HTML}{5D6470}
\definecolor{Ink}{HTML}{20262D}
\definecolor{Muted}{HTML}{66717B}
\definecolor{Panel}{HTML}{F5F7F8}
\definecolor{RuleLine}{HTML}{D7DEE3}

\hypersetup{
  pdftitle={GitHub Advanced Security Supported-Language Secure Coding Standards Catalog},
  pdfauthor={Application Security Engineering},
  colorlinks=true,
  linkcolor=Navy,
  urlcolor=Teal,
  citecolor=Navy
}

\pagestyle{fancy}
\fancyhf{}
\lhead{\textcolor{Muted}{GHAS Secure Coding Standards}}
\rhead{\textcolor{Muted}{Enterprise Baseline v1.0}}
\cfoot{\textcolor{Muted}{\thepage\ / \pageref*{LastPage}}}
\renewcommand{\headrulewidth}{0.3pt}
\renewcommand{\footrulewidth}{0pt}

\titleformat{\section}{\Large\bfseries\color{Navy}}{\thesection}{0.65em}{}
\titleformat{\subsection}{\large\bfseries\color{Teal}}{\thesubsection}{0.6em}{}
\titleformat{\subsubsection}{\normalsize\bfseries\color{Ink}}{\thesubsubsection}{0.5em}{}
\titlespacing*{\section}{0pt}{1.7em}{0.7em}
\titlespacing*{\subsection}{0pt}{1.3em}{0.5em}
\setcounter{tocdepth}{2}
\setcounter{secnumdepth}{2}

\newcommand{\code}[1]{\texttt{\small #1}}
\newcommand{\prioritybadge}[2]{\tcbox[on line,boxrule=0pt,arc=1.2mm,boxsep=0.7mm,left=1mm,right=1mm,top=0.4mm,bottom=0.4mm,colback=#2!12,coltext=#2]{\scriptsize\bfseries #1}}
\newcommand{\pill}[1]{\tcbox[on line,boxrule=0pt,arc=1mm,boxsep=0.6mm,left=0.8mm,right=0.8mm,top=0.3mm,bottom=0.3mm,colback=LightTeal,coltext=Teal]{\scriptsize\bfseries #1}}

\newtcolorbox{rulebox}[1][]{
  enhanced, breakable, colback=white, colframe=RuleLine,
  boxrule=0.55pt, arc=1.4mm, left=2.5mm,right=2.5mm,top=1.5mm,bottom=1.5mm,
  borderline west={2.2pt}{0pt}{Teal},
  fonttitle=\bfseries\color{Navy}, coltitle=Navy,
  colbacktitle=Panel, attach boxed title to top left={xshift=1.5mm,yshift*=-1mm},
  boxed title style={boxrule=0pt,arc=1mm,colback=Panel},
  before skip=8pt, after skip=8pt, #1
}

\newtcolorbox{callout}[2][]{
  enhanced, breakable, colback=#2!5, colframe=#2!40,
  boxrule=0.6pt, arc=1.3mm, left=3mm,right=3mm,top=2mm,bottom=2mm,
  #1
}

\lstdefinestyle{securecode}{
  basicstyle=\ttfamily\footnotesize,
  columns=fullflexible,
  keepspaces=true,
  showstringspaces=false,
  breaklines=true,
  breakatwhitespace=false,
  frame=single,
  rulecolor=\color{RuleLine},
  backgroundcolor=\color{Panel},
  xleftmargin=2mm,
  xrightmargin=1mm,
  aboveskip=6pt,
  belowskip=8pt,
  tabsize=2
}
\lstset{style=securecode}

\newcolumntype{Y}{>{\RaggedRight\arraybackslash}X}
\newcolumntype{P}[1]{>{\RaggedRight\arraybackslash}p{#1}}
\renewcommand{\arraystretch}{1.15}
\setlength{\emergencystretch}{3em}


\begin{document}
\hypersetup{pdftitle={Enterprise Secure Coding Standard --- JavaScript},pdfauthor={Application Security Engineering}}
\fancyhead[L]{\textcolor{Muted}{Enterprise Secure Coding Standard --- JavaScript}}
\fancyhead[R]{\textcolor{Muted}{v1.0 --- 23 September 2026}}
\hypersetup{pageanchor=false}
\begin{titlepage}
\thispagestyle{empty}
\vspace*{0.85in}
{\color{Teal}\rule{1.2in}{3pt}}\\[0.32in]
{\Huge\bfseries\color{Navy}Enterprise Secure Coding Standard --- JavaScript\par}
\vspace{0.18in}
{\Large\color{Teal}GitHub Advanced Security / CodeQL Enforcement Profile\par}
\vspace{0.55in}
\begin{callout}{Teal}
\textbf{Purpose.} Provide a standalone, traceable secure-coding standard that connects normative implementation rules to CWE, CodeQL detection, companion verification, and SSDLC enforcement.
\end{callout}
\vfill
\begin{tabular}{@{}ll@{}}
\textbf{Standard version} & 1.0 \\
\textbf{Verified} & 23 September 2026 \\
\textbf{Document status} & Enterprise baseline \\
\textbf{Scope} & JavaScript application and library code \\
\textbf{Source policy} & Official and authoritative primary sources only \\
\textbf{Normative terms} & MUST, MUST NOT, SHOULD, SHOULD NOT, MAY \\
\end{tabular}
\vspace{0.55in}
\end{titlepage}
\hypersetup{pageanchor=true}
\pagenumbering{roman}
\tableofcontents
\clearpage
\pagenumbering{arabic}

\section{Executive Summary}
This standard defines the secure-coding baseline for this language and combines applicable cross-language controls with language-specific implementation rules, CodeQL detection, companion automation, and SSDLC gates. It is intended for developers, security engineers, security champions, reviewers, DevSecOps, vulnerability-management teams, GHAS administrators, architects, engineering managers, and assurance functions.

\begin{callout}{Navy}
\textbf{Interpretation boundary.} A clean CodeQL result does not demonstrate that an application is secure. CodeQL is one verification control inside a broader SSDLC that also requires secure design, runtime validation, authentication/authorization testing, dependency and secret controls, dynamic testing, configuration review, and operational monitoring.
\end{callout}

\textbf{Normative interpretation.} MUST/MUST NOT are mandatory unless an approved time-bounded exception exists. SHOULD/SHOULD NOT require documented engineering rationale when deviated from. MAY denotes an optional technique.

\section{Language Security Profile}
JavaScript is dynamically typed and strongly runtime-context dependent. Browser DOM APIs, Node.js/server APIs, prototypes, dynamic execution, process APIs, templating, asynchronous resource behavior, and package supply-chain semantics dominate the security model.
\section{Authority and Standards Baseline}
\begin{callout}{Teal}
\textbf{Classification:} Derived Secure Coding Baseline\\
\textbf{Primary authority:} Node.js security/API documentation + ECMAScript/runtime semantics\\
\textbf{Supporting authorities:} current GitHub CodeQL documentation, NIST SSDF, MITRE CWE, OWASP ASVS where applicable, and the platform-specific sources listed in the reference catalog.
\end{callout}
This classification describes the provenance of the standard; it does not imply that a single external publication completely covers the enterprise control set. Where the language lacks one comprehensive secure-coding standard, this document derives a baseline only from official/authoritative first-party sources and preserves the canonical enterprise rules from the supplied catalog.

\section{CodeQL Support Profile}
\begin{tabularx}{\textwidth}{@{}P{1.9in}Y@{}}
\toprule
\textbf{Field} & \textbf{Value}\\
\midrule
Language & JavaScript\\
CodeQL identifier & \code{javascript-typescript}\\
Supported versions & ECMAScript 2022 or lower\\
Extraction / build mode & none\\
Build required & No\\
Key documented frameworks/libraries & Browser, Node, Express, Fastify, React, Vue, Angular, common DB/template/network libraries\\
Current limitations & Browser and server-side JavaScript have different runtime trust boundaries.\\
Official GitHub documentation & \url{https://codeql.github.com/docs/codeql-overview/supported-languages-and-frameworks/}\\
Verification date & 23 September 2026\\
\bottomrule
\end{tabularx}
\begin{callout}{Teal}
\textbf{Maintenance note.} Current GitHub documentation was reverified on 23 September 2026. No material language/version change requiring a canonical rule-ID change was identified during this generation pass.
\end{callout}
\section{Security Domain Applicability}
\small
\setlength{\tabcolsep}{3pt}
\begin{longtable}{@{}P{1.8in}P{0.85in}P{3.75in}@{}}
\toprule
\textbf{Domain} & \textbf{Applicability} & \textbf{Notes}\\
\midrule
\endfirsthead
\toprule
\textbf{Domain} & \textbf{Applicability} & \textbf{Notes}\\
\midrule
\endhead
Input Validation & Required & Canonical cross-language baseline applies.\\
Output Encoding & Required & Canonical cross-language baseline applies.\\
Injection Prevention & Required & Canonical cross-language baseline applies.\\
Memory Safety & Scoped & Runtime-managed; conditional for native add-ons\\
Authentication & Required & Canonical cross-language baseline applies.\\
Authorization & Required & Canonical cross-language baseline applies.\\
Cryptography & Required & Canonical cross-language baseline applies.\\
Secrets & Required & Canonical cross-language baseline applies.\\
Error / Exception Handling & Required & Canonical cross-language baseline applies.\\
Logging / Monitoring & Required & Canonical cross-language baseline applies.\\
File / Path Security & Required & Canonical cross-language baseline applies.\\
Serialization / Deserialization & Required & Canonical cross-language baseline applies.\\
Network Communications & Required & Canonical cross-language baseline applies.\\
Web Application Security & Required & Required for browser/web/server contexts\\
Dependency / Package Security & Required & Canonical cross-language baseline applies.\\
Resource Management & Required & Canonical cross-language baseline applies.\\
Concurrency / Parallelism & Required & Canonical cross-language baseline applies.\\
Dynamic Code Execution & Required & Canonical cross-language baseline applies.\\
Native / Foreign Interfaces & Conditional & Conditional: native add-ons / FFI\\
Compiler / Build Security & Required & Canonical cross-language baseline applies.\\
\bottomrule
\end{longtable}
\normalsize
All domains are evaluated independently of CodeQL coverage. A domain can remain a normative requirement even when the primary verification control is dynamic testing, a platform tool, another GitHub security capability, or manual review.
\section{Applicable Cross-Language Secure Coding Requirements}
The following canonical `SCS-ALL-*` controls are copied verbatim from the governing catalog, with only explicit prevention/verification labels added. Conditional scope is governed by the applicability table above.
\begin{rulebox}[title={\texttt{SCS-ALL-INPUT-001} \hfill \prioritybadge{High Priority}{High}}]
\textbf{Domain.} Input Validation

\textbf{Requirement.} All data crossing a trust boundary MUST be validated for expected type, syntax, length, range, structure, and business constraints before security-sensitive use.

\textbf{Rationale.} Type systems and parser success do not establish that attacker-controlled data is authorized or semantically safe.

\textbf{Risk.} Injection, parser abuse, resource exhaustion, authorization bypass

\begin{description}[leftmargin=2.65cm,labelwidth=2.3cm,style=nextline,itemsep=1pt,topsep=3pt]
  \item[CWE] CWE-20
  \item[OWASP] OWASP ASVS 5.0 --- validation/business logic
  \item[CERT] Not mapped
  \item[CodeQL] \textbf{Partial}; Language-specific taint/data-flow queries
  \item[Enforcement] CodeQL + schema/parser validation + unit/fuzz tests
  \item[Prevention] Schema/domain validation at trust boundaries before security-sensitive use.
  \item[Verification] CodeQL + schema/parser validation + unit/fuzz tests
  \item[Automation] Partial
  \item[Sources] \texttt{NIST-SSDF, OWASP-ASVS, MITRE-CWE, JAVA-SCG}
\end{description}
\end{rulebox}
\begin{rulebox}[title={\texttt{SCS-ALL-OUTPUT-001} \hfill \prioritybadge{High Priority}{High}}]
\textbf{Domain.} Output Encoding

\textbf{Requirement.} Untrusted data MUST be encoded or escaped for the exact output context; developers MUST prefer framework APIs that perform contextual escaping by default.

\textbf{Rationale.} Encoding is context-sensitive; HTML, JavaScript, URLs, headers, XML, SQL, and shell contexts are not interchangeable.

\textbf{Risk.} XSS, response splitting, code/command interpretation

\begin{description}[leftmargin=2.65cm,labelwidth=2.3cm,style=nextline,itemsep=1pt,topsep=3pt]
  \item[CWE] CWE-79, CWE-116
  \item[OWASP] OWASP ASVS 5.0 --- encoding/sanitization
  \item[CERT] Not mapped
  \item[CodeQL] \textbf{Partial}; Language/framework-specific XSS and response-splitting queries
  \item[Enforcement] CodeQL + framework-safe templates + tests
  \item[Prevention] Context-aware encoders or framework escaping APIs.
  \item[Verification] CodeQL + framework-safe templates + tests
  \item[Automation] Partial
  \item[Sources] \texttt{OWASP-ASVS, MITRE-CWE, GH-BUILTIN}
\end{description}
\end{rulebox}
\begin{rulebox}[title={\texttt{SCS-ALL-INJ-001} \hfill \prioritybadge{Critical Control}{Critical}}]
\textbf{Domain.} Injection Prevention

\textbf{Requirement.} Untrusted data MUST NOT be concatenated into SQL, OS commands, code, templates, expressions, or other interpreters; parameterized or structured APIs MUST be used where available.

\textbf{Rationale.} Separating code from data removes attacker control over interpreter syntax.

\textbf{Risk.} SQL/NoSQL injection, command injection, code injection, template/expression injection

\begin{description}[leftmargin=2.65cm,labelwidth=2.3cm,style=nextline,itemsep=1pt,topsep=3pt]
  \item[CWE] CWE-78, CWE-89, CWE-94, CWE-95, CWE-943
  \item[OWASP] OWASP ASVS 5.0 --- injection prevention
  \item[CERT] Not mapped
  \item[CodeQL] \textbf{Direct}; Representative: cpp/sql-injection; cs/sql-injection; go/sql-injection; java/sql-injection; js/sql-injection; py/sql-injection; rb/sql-injection; rust/sql-injection; swift/sql-injection
  \item[Enforcement] CodeQL required check + API-level parameterization
  \item[Prevention] Structured/parameterized APIs; keep untrusted data out of interpreter syntax.
  \item[Verification] CodeQL required check + API-level parameterization
  \item[Automation] Yes
  \item[Sources] \texttt{GH-BUILTIN, MITRE-CWE, OWASP-ASVS}
\end{description}
\end{rulebox}
\begin{rulebox}[title={\texttt{SCS-ALL-MEM-001} \hfill \prioritybadge{Critical Control}{Critical}}]
\textbf{Domain.} Memory Safety

\textbf{Requirement.} Code that can perform unchecked memory access MUST establish bounds, initialization, lifetime, ownership, and aliasing invariants before access; memory-safe language code SHOULD remain inside safe subsets.

\textbf{Rationale.} Memory corruption can become arbitrary code execution or information disclosure.

\textbf{Risk.} Out-of-bounds access, use-after-free, double free, invalid pointer use

\begin{description}[leftmargin=2.65cm,labelwidth=2.3cm,style=nextline,itemsep=1pt,topsep=3pt]
  \item[CWE] CWE-119, CWE-416, CWE-415, CWE-825
  \item[OWASP] Context-dependent / not forced
  \item[CERT] SEI CERT C/C++ memory rules
  \item[CodeQL] \textbf{Partial}; Representative: cpp/use-after-free; rust/access-after-lifetime-ended
  \item[Enforcement] CodeQL + compiler/runtime checks + sanitizers/dynamic tests + review of unsafe/native boundaries
  \item[Prevention] Safe abstractions plus explicit bounds/lifetime/ownership invariants at unsafe boundaries.
  \item[Verification] CodeQL + compiler/runtime checks + sanitizers/dynamic tests + review of unsafe/native boundaries
  \item[Automation] Partial
  \item[Sources] \texttt{CERT-C, CERT-CPP, RUST-UNSAFE, SWIFT-MEM, GH-BUILTIN}
\end{description}
\end{rulebox}
\begin{rulebox}[title={\texttt{SCS-ALL-AUTHN-001} \hfill \prioritybadge{Critical Control}{Critical}}]
\textbf{Domain.} Authentication

\textbf{Requirement.} Applications MUST use a vetted platform/framework/identity-provider authentication mechanism and MUST NOT implement custom credential cryptography or authentication protocols without specialized review.

\textbf{Rationale.} Authentication implementations are protocol- and state-sensitive and commonly fail at credential, session, recovery, or trust boundaries.

\textbf{Risk.} Authentication bypass, credential compromise

\begin{description}[leftmargin=2.65cm,labelwidth=2.3cm,style=nextline,itemsep=1pt,topsep=3pt]
  \item[CWE] CWE-287, CWE-306
  \item[OWASP] OWASP ASVS 5.0 --- authentication
  \item[CERT] Not mapped
  \item[CodeQL] \textbf{Partial}; Framework-specific authentication queries
  \item[Enforcement] Architecture review + framework configuration tests + DAST/IAST where applicable
  \item[Prevention] Vetted framework/identity-provider controls; avoid custom authentication cryptography/protocols.
  \item[Verification] Architecture review + framework configuration tests + DAST/IAST where applicable
  \item[Automation] Partial
  \item[Sources] \texttt{OWASP-ASVS, NIST-SSDF, MITRE-CWE}
\end{description}
\end{rulebox}
\begin{rulebox}[title={\texttt{SCS-ALL-AUTHZ-001} \hfill \prioritybadge{Critical Control}{Critical}}]
\textbf{Domain.} Authorization

\textbf{Requirement.} Authorization MUST be enforced server-side at each protected operation and object boundary, SHOULD default deny, and MUST be based on trusted identity/authorization state rather than client-supplied claims alone.

\textbf{Rationale.} Authenticated users can still request objects or actions they are not permitted to access.

\textbf{Risk.} Missing authorization, IDOR/BOLA, privilege escalation

\begin{description}[leftmargin=2.65cm,labelwidth=2.3cm,style=nextline,itemsep=1pt,topsep=3pt]
  \item[CWE] CWE-285, CWE-862, CWE-863, CWE-639
  \item[OWASP] OWASP ASVS 5.0 --- access control
  \item[CERT] Not mapped
  \item[CodeQL] \textbf{Partial}; Language/framework-specific missing authorization / IDOR queries
  \item[Enforcement] CodeQL where modeled + authorization tests + architecture/code review
  \item[Prevention] Server-side, default-deny checks at protected object and operation boundaries.
  \item[Verification] CodeQL where modeled + authorization tests + architecture/code review
  \item[Automation] Partial
  \item[Sources] \texttt{OWASP-ASVS, MITRE-CWE, GH-BUILTIN}
\end{description}
\end{rulebox}
\begin{rulebox}[title={\texttt{SCS-ALL-CRYPTO-001} \hfill \prioritybadge{Critical Control}{Critical}}]
\textbf{Domain.} Cryptography

\textbf{Requirement.} Applications MUST use current platform cryptographic APIs and cryptographically secure random generators for security-sensitive values; custom cryptographic primitives and obsolete/broken algorithms MUST NOT be used.

\textbf{Rationale.} Cryptographic strength depends on algorithm, mode, parameters, randomness, key management, and API usage.

\textbf{Risk.} Predictable tokens, weak encryption/signatures, key compromise

\begin{description}[leftmargin=2.65cm,labelwidth=2.3cm,style=nextline,itemsep=1pt,topsep=3pt]
  \item[CWE] CWE-327, CWE-326, CWE-338
  \item[OWASP] OWASP ASVS 5.0 --- cryptography
  \item[CERT] Not mapped
  \item[CodeQL] \textbf{Partial}; Weak-crypto / insecure-randomness / TLS queries vary by language
  \item[Enforcement] CodeQL + platform API policy + cryptographic design review
  \item[Prevention] Current platform cryptographic APIs and approved key-management practices.
  \item[Verification] CodeQL + platform API policy + cryptographic design review
  \item[Automation] Partial
  \item[Sources] \texttt{OWASP-ASVS, MITRE-CWE, DOTNET-CRYPTO, GO-RAND, JAVA-26, APPLE-CRYPTO}
\end{description}
\end{rulebox}
\begin{rulebox}[title={\texttt{SCS-ALL-SECRET-001} \hfill \prioritybadge{Critical Control}{Critical}}]
\textbf{Domain.} Secrets

\textbf{Requirement.} Credentials, API tokens, private keys, and other secrets MUST NOT be hard-coded in source or written to logs; repositories SHOULD enable GitHub secret scanning and push protection for supported secret types.

\textbf{Rationale.} Source history and logs are durable distribution channels; exposed credentials can be replayed.

\textbf{Risk.} Credential exposure and unauthorized access

\begin{description}[leftmargin=2.65cm,labelwidth=2.3cm,style=nextline,itemsep=1pt,topsep=3pt]
  \item[CWE] CWE-798, CWE-532, CWE-312
  \item[OWASP] OWASP ASVS 5.0 --- secrets/data protection
  \item[CERT] Not mapped
  \item[CodeQL] \textbf{Requires another control}; GitHub secret scanning / push protection rather than CodeQL for primary enforcement
  \item[Enforcement] Secret scanning + push protection + approved secret store + rotation process
  \item[Prevention] External secret stores, least privilege, and rotation/revocation after exposure.
  \item[Verification] Secret scanning + push protection + approved secret store + rotation process
  \item[Automation] Yes
  \item[Sources] \texttt{GH-SECRET, GH-PUSH, MITRE-CWE, NIST-SSDF}
\end{description}
\end{rulebox}
\begin{rulebox}[title={\texttt{SCS-ALL-ERROR-001} \hfill \prioritybadge{High Priority}{High}}]
\textbf{Domain.} Error / Exception Handling

\textbf{Requirement.} Security failures MUST fail closed; externally visible errors MUST NOT disclose secrets, credentials, stack traces, filesystem paths, internal queries, or other sensitive implementation details.

\textbf{Rationale.} Error handling is part of the security boundary and can otherwise disclose exploitable state or bypass checks.

\textbf{Risk.} Information exposure, fail-open authorization/authentication behavior

\begin{description}[leftmargin=2.65cm,labelwidth=2.3cm,style=nextline,itemsep=1pt,topsep=3pt]
  \item[CWE] CWE-209, CWE-703
  \item[OWASP] OWASP ASVS 5.0 --- error handling/logging
  \item[CERT] Not mapped
  \item[CodeQL] \textbf{Partial}; Information-disclosure and unchecked-error queries vary
  \item[Enforcement] Code review + negative-path tests + logging policy + CodeQL where available
  \item[Prevention] Fail closed and avoid sensitive detail in externally visible errors.
  \item[Verification] Code review + negative-path tests + logging policy + CodeQL where available
  \item[Automation] Partial
  \item[Sources] \texttt{OWASP-ASVS, MITRE-CWE, NIST-SSDF}
\end{description}
\end{rulebox}
\begin{rulebox}[title={\texttt{SCS-ALL-LOG-001} \hfill \prioritybadge{High Priority}{High}}]
\textbf{Domain.} Logging / Monitoring

\textbf{Requirement.} Security-relevant events SHOULD be logged with structured fields; untrusted data MUST be encoded for the logging sink and secrets/sensitive payloads MUST NOT be logged.

\textbf{Rationale.} Logs are both security evidence and attacker-controlled sinks; injection or sensitive data exposure weakens monitoring and confidentiality.

\textbf{Risk.} Log injection, sensitive data exposure, audit gaps

\begin{description}[leftmargin=2.65cm,labelwidth=2.3cm,style=nextline,itemsep=1pt,topsep=3pt]
  \item[CWE] CWE-117, CWE-532
  \item[OWASP] OWASP ASVS 5.0 --- logging
  \item[CERT] Not mapped
  \item[CodeQL] \textbf{Partial}; Log-injection / cleartext-logging queries exist for multiple languages
  \item[Enforcement] CodeQL + centralized logging policy + tests
  \item[Prevention] Structured security logging without secrets; sanitize/encode attacker-controlled fields.
  \item[Verification] CodeQL + centralized logging policy + tests
  \item[Automation] Partial
  \item[Sources] \texttt{OWASP-ASVS, MITRE-CWE, GH-BUILTIN}
\end{description}
\end{rulebox}
\begin{rulebox}[title={\texttt{SCS-ALL-FILE-001} \hfill \prioritybadge{High Priority}{High}}]
\textbf{Domain.} File / Path Security

\textbf{Requirement.} File operations MUST constrain attacker-controlled paths to an intended root, avoid unsafe symbolic-link/TOCTOU assumptions, use secure temporary-file APIs, and validate archive extraction destinations.

\textbf{Rationale.} Path normalization alone is not an authorization decision and filesystem state can change between checks and use.

\textbf{Risk.} Path traversal, arbitrary file read/write, symlink race, unsafe temp files, Zip Slip

\begin{description}[leftmargin=2.65cm,labelwidth=2.3cm,style=nextline,itemsep=1pt,topsep=3pt]
  \item[CWE] CWE-22, CWE-59, CWE-367, CWE-377
  \item[OWASP] OWASP ASVS 5.0 --- file handling
  \item[CERT] Not mapped
  \item[CodeQL] \textbf{Direct/Partial}; Path-injection/traversal, archive extraction, temporary-file, and TOCTOU queries vary by language
  \item[Enforcement] CodeQL + filesystem permission tests + archive policy
  \item[Prevention] Root confinement, safe temporary/archive APIs, and handle-based authorization where possible.
  \item[Verification] CodeQL + filesystem permission tests + archive policy
  \item[Automation] Partial
  \item[Sources] \texttt{GH-BUILTIN, MITRE-CWE, PY-TAR, NIST-SSDF}
\end{description}
\end{rulebox}
\begin{rulebox}[title={\texttt{SCS-ALL-DESER-001} \hfill \prioritybadge{Critical Control}{Critical}}]
\textbf{Domain.} Serialization / Deserialization

\textbf{Requirement.} Untrusted input MUST NOT be deserialized by APIs that can instantiate arbitrary types or invoke attacker-influenced code; data-only formats and restrictive/safe loaders MUST be preferred.

\textbf{Rationale.} Powerful object deserializers can convert data parsing into code execution.

\textbf{Risk.} Remote code execution, object injection, denial of service

\begin{description}[leftmargin=2.65cm,labelwidth=2.3cm,style=nextline,itemsep=1pt,topsep=3pt]
  \item[CWE] CWE-502
  \item[OWASP] OWASP ASVS 5.0 --- deserialization
  \item[CERT] Not mapped
  \item[CodeQL] \textbf{Direct/Partial}; Representative: cs/unsafe-deserialization-untrusted-input; java/unsafe-deserialization; py/unsafe-deserialization; rb/unsafe-deserialization
  \item[Enforcement] CodeQL + API allowlist + schema validation
  \item[Prevention] Data-only formats or restrictive safe loaders; avoid arbitrary-object deserialization.
  \item[Verification] CodeQL + API allowlist + schema validation
  \item[Automation] Yes/Partial
  \item[Sources] \texttt{DOTNET-BF, PY-PICKLE, RUBY-MARSHAL, JAVA-SCG, GH-BUILTIN}
\end{description}
\end{rulebox}
\begin{rulebox}[title={\texttt{SCS-ALL-NET-001} \hfill \prioritybadge{Critical Control}{Critical}}]
\textbf{Domain.} Network Communications

\textbf{Requirement.} Sensitive or authenticated network traffic MUST use TLS with normal certificate and hostname verification; production code MUST NOT disable or bypass trust validation except through a documented, narrowly scoped exception.

\textbf{Rationale.} Encryption without peer authentication remains vulnerable to interception and impersonation.

\textbf{Risk.} MITM, cleartext disclosure, downgrade

\begin{description}[leftmargin=2.65cm,labelwidth=2.3cm,style=nextline,itemsep=1pt,topsep=3pt]
  \item[CWE] CWE-295, CWE-297, CWE-319, CWE-757
  \item[OWASP] OWASP ASVS 5.0 --- communications
  \item[CERT] Not mapped
  \item[CodeQL] \textbf{Direct/Partial}; Representative: go/disabled-certificate-check; swift/insecure-tls; TLS/certificate queries in other packs
  \item[Enforcement] CodeQL + platform network policy + integration tests
  \item[Prevention] TLS with certificate and hostname verification; constrained outbound destinations where needed.
  \item[Verification] CodeQL + platform network policy + integration tests
  \item[Automation] Partial
  \item[Sources] \texttt{OWASP-ASVS, MITRE-CWE, APPLE-ATS, GH-BUILTIN}
\end{description}
\end{rulebox}
\begin{rulebox}[title={\texttt{SCS-ALL-WEB-001} \hfill \prioritybadge{Critical Control}{Critical}}]
\textbf{Domain.} Web Application Security

\textbf{Requirement.} Web applications MUST use framework protections for contextual output encoding, anti-CSRF where browser credentials are ambient, secure cookies/session handling, constrained CORS, validated redirects, and outbound-request controls against SSRF.

\textbf{Rationale.} Web security depends on browser/proxy semantics and trust boundaries not provided by the base programming language.

\textbf{Risk.} XSS, CSRF, SSRF, open redirect, session compromise

\begin{description}[leftmargin=2.65cm,labelwidth=2.3cm,style=nextline,itemsep=1pt,topsep=3pt]
  \item[CWE] CWE-79, CWE-352, CWE-918, CWE-601
  \item[OWASP] OWASP ASVS 5.0 --- web/browser/session controls
  \item[CERT] Not mapped
  \item[CodeQL] \textbf{Direct/Partial}; Representative: go/request-forgery; js/request-forgery; rust/request-forgery; framework-specific XSS/CSRF queries
  \item[Enforcement] CodeQL + framework config tests + DAST/IAST + architecture review
  \item[Prevention] Framework browser/session protections plus server-side validation and authorization.
  \item[Verification] CodeQL + framework config tests + DAST/IAST + architecture review
  \item[Automation] Partial
  \item[Sources] \texttt{OWASP-ASVS, GH-BUILTIN, ASPNET-CSRF, RAILS-SEC}
\end{description}
\end{rulebox}
\begin{rulebox}[title={\texttt{SCS-ALL-DEP-001} \hfill \prioritybadge{Critical Control}{Critical}}]
\textbf{Domain.} Dependency / Package Security

\textbf{Requirement.} Dependencies MUST come from approved sources, SHOULD be locked/pinned reproducibly, MUST be reviewed for known vulnerabilities, and MUST be updated or mitigated according to risk; dependency-review gates SHOULD protect pull requests where available.

\textbf{Rationale.} A secure application can be compromised by a vulnerable or malicious transitive dependency.

\textbf{Risk.} Known-vulnerable component, dependency confusion, malicious package/action

\begin{description}[leftmargin=2.65cm,labelwidth=2.3cm,style=nextline,itemsep=1pt,topsep=3pt]
  \item[CWE] CWE-829
  \item[OWASP] OWASP ASVS 5.0 --- dependency/supply-chain controls
  \item[CERT] Not mapped
  \item[CodeQL] \textbf{Requires another control}; Dependabot alerts/security updates + dependency review
  \item[Enforcement] Dependabot + dependency review + lockfiles + approved registries
  \item[Prevention] Approved registries, reproducible dependency state, dependency review, and timely remediation.
  \item[Verification] Dependabot + dependency review + lockfiles + approved registries
  \item[Automation] Yes
  \item[Sources] \texttt{GH-DEP-REVIEW, GH-DEPENDABOT, NIST-SSDF, JAVA-SCG}
\end{description}
\end{rulebox}
\begin{rulebox}[title={\texttt{SCS-ALL-RESOURCE-001} \hfill \prioritybadge{High Priority}{High}}]
\textbf{Domain.} Resource Management

\textbf{Requirement.} Untrusted inputs MUST be bounded before allocation, decompression, regex evaluation, recursion, parsing, fan-out, or other resource-amplifying work; services SHOULD enforce timeouts and quotas.

\textbf{Rationale.} Small attacker requests can trigger disproportionate CPU, memory, storage, connection, or thread consumption.

\textbf{Risk.} Denial of service, algorithmic complexity abuse

\begin{description}[leftmargin=2.65cm,labelwidth=2.3cm,style=nextline,itemsep=1pt,topsep=3pt]
  \item[CWE] CWE-400, CWE-770, CWE-789, CWE-1333
  \item[OWASP] OWASP ASVS 5.0 --- availability/resource controls
  \item[CERT] Not mapped
  \item[CodeQL] \textbf{Partial}; Allocation-size, regex, decompression and other resource queries vary by language
  \item[Enforcement] CodeQL + fuzz/load tests + runtime quotas/timeouts
  \item[Prevention] Explicit sizes, quotas, timeouts, backpressure, and complexity limits.
  \item[Verification] CodeQL + fuzz/load tests + runtime quotas/timeouts
  \item[Automation] Partial
  \item[Sources] \texttt{JAVA-SCG, GO-SEC, NODE-SEC, MITRE-CWE}
\end{description}
\end{rulebox}
\begin{rulebox}[title={\texttt{SCS-ALL-CONC-001} \hfill \prioritybadge{High Priority}{High}}]
\textbf{Domain.} Concurrency / Parallelism

\textbf{Requirement.} Shared mutable state MUST use language-appropriate synchronization/ownership mechanisms; security decisions MUST NOT rely on filesystem or shared state remaining unchanged between separate check and use operations.

\textbf{Rationale.} Races can invalidate authorization, resource, memory, or filesystem assumptions.

\textbf{Risk.} Race conditions, TOCTOU, corrupted shared state

\begin{description}[leftmargin=2.65cm,labelwidth=2.3cm,style=nextline,itemsep=1pt,topsep=3pt]
  \item[CWE] CWE-362, CWE-367
  \item[OWASP] Context-dependent / not forced
  \item[CERT] SEI CERT concurrency rules where applicable
  \item[CodeQL] \textbf{Partial}; TOCTOU/static queries where available; dynamic race detectors for Go and native code
  \item[Enforcement] CodeQL + race detector/thread sanitizer where supported + concurrency tests
  \item[Prevention] Language-appropriate synchronization/ownership and race-aware testing.
  \item[Verification] CodeQL + race detector/thread sanitizer where supported + concurrency tests
  \item[Automation] Partial
  \item[Sources] \texttt{GO-SEC, CERT-C, CERT-CPP, MITRE-CWE}
\end{description}
\end{rulebox}
\begin{rulebox}[title={\texttt{SCS-ALL-DYN-001} \hfill \prioritybadge{Critical Control}{Critical}}]
\textbf{Domain.} Dynamic Code Execution

\textbf{Requirement.} Untrusted data MUST NOT reach eval-like APIs, runtime code generation, dynamic loaders, reflective invocation targets, or interpreters unless constrained by a purpose-built parser and explicit allowlist.

\textbf{Rationale.} Dynamic execution turns attacker-controlled data into executable behavior.

\textbf{Risk.} Code injection, arbitrary execution

\begin{description}[leftmargin=2.65cm,labelwidth=2.3cm,style=nextline,itemsep=1pt,topsep=3pt]
  \item[CWE] CWE-94, CWE-95
  \item[OWASP] OWASP ASVS 5.0 --- injection
  \item[CERT] Not mapped
  \item[CodeQL] \textbf{Direct/Partial}; Code-injection / dynamic-code queries vary by language
  \item[Enforcement] CodeQL + code review + API restrictions
  \item[Prevention] Avoid eval-like sinks; use parsers and explicit allowlists if dynamic behavior is unavoidable.
  \item[Verification] CodeQL + code review + API restrictions
  \item[Automation] Partial
  \item[Sources] \texttt{GH-BUILTIN, MITRE-CWE, OWASP-ASVS}
\end{description}
\end{rulebox}
\begin{rulebox}[title={\texttt{SCS-ALL-FFI-001} \hfill \prioritybadge{Critical Control}{Critical}}]
\textbf{Domain.} Native / Foreign Interfaces

\textbf{Requirement.} FFI/native boundaries MUST be treated as unsafe trust boundaries: signatures, ownership, lifetime, nullability, buffer lengths, thread-safety, and error conventions MUST be validated and documented; a safe wrapper SHOULD isolate raw calls.

\textbf{Rationale.} Foreign runtimes can invalidate guarantees made by a memory-safe or type-safe language.

\textbf{Risk.} Memory corruption, lifetime violations, ABI mismatch

\begin{description}[leftmargin=2.65cm,labelwidth=2.3cm,style=nextline,itemsep=1pt,topsep=3pt]
  \item[CWE] CWE-119, CWE-416, CWE-704
  \item[OWASP] Context-dependent / not forced
  \item[CERT] SEI CERT / language-specific FFI guidance
  \item[CodeQL] \textbf{Partial}; Native-side CodeQL where source is analyzed; limited cross-language guarantees
  \item[Enforcement] Boundary review + native/static analysis + sanitizers + tests
  \item[Prevention] Small audited wrappers with documented pointer, lifetime, ownership, and ABI contracts.
  \item[Verification] Boundary review + native/static analysis + sanitizers + tests
  \item[Automation] Partial
  \item[Sources] \texttt{RUST-FFI, SWIFT-INTEROP, JAVA-SCG, DOTNET-SEC}
\end{description}
\end{rulebox}
\begin{rulebox}[title={\texttt{SCS-ALL-BUILD-001} \hfill \prioritybadge{High Priority}{High}}]
\textbf{Domain.} Compiler / Build Security

\textbf{Requirement.} Builds MUST use supported toolchains, fail on security-relevant static-analysis findings according to policy, preserve reproducible dependency state where the ecosystem supports lockfiles, and SHOULD run language-native security diagnostics/tests in CI.

\textbf{Rationale.} Security controls that are not reproducible and enforced in CI degrade over time and across developer environments.

\textbf{Risk.} Unreviewed defects, vulnerable dependencies, inconsistent builds

\begin{description}[leftmargin=2.65cm,labelwidth=2.3cm,style=nextline,itemsep=1pt,topsep=3pt]
  \item[CWE] CWE-1104 (where applicable)
  \item[OWASP] Context-dependent / not forced
  \item[CERT] Not mapped
  \item[CodeQL] \textbf{Requires another control}; CodeQL check runs plus language-native compiler/linter/test/dependency checks
  \item[Enforcement] CI/CD policy + CodeQL + ecosystem-native tooling
  \item[Prevention] Supported toolchains, reproducible builds, static analysis, and security diagnostics in CI.
  \item[Verification] CI/CD policy + CodeQL + ecosystem-native tooling
  \item[Automation] Yes/Partial
  \item[Sources] \texttt{NIST-SSDF, GH-LANG, GH-DEP-REVIEW, GO-SEC, TS-STRICT}
\end{description}
\end{rulebox}
\section{Language-Specific Secure Coding Requirements}
\begin{rulebox}[title={\texttt{SCS-JS-DYN-001} \hfill \prioritybadge{Critical Control}{Critical}}]
\textbf{Domain.} Dynamic Code Execution

\textbf{Requirement.} Untrusted strings MUST NOT be passed to eval, Function constructors, dynamic code loaders, template compilation, or equivalent executable sinks.

\textbf{Rationale.} JavaScript's dynamic execution features can interpret attacker input directly as code.

\textbf{Risk.} Code injection

\begin{description}[leftmargin=2.65cm,labelwidth=2.3cm,style=nextline,itemsep=1pt,topsep=3pt]
  \item[CWE] CWE-94, CWE-95
  \item[OWASP] OWASP ASVS 5.0 --- injection
  \item[CERT] Not mapped
  \item[CodeQL] \textbf{Direct}; JavaScript code-injection queries
  \item[Enforcement] CodeQL + lint/API restrictions
  \item[Prevention] Avoid eval-like sinks; use parsers and explicit allowlists if dynamic behavior is unavoidable.
  \item[Verification] CodeQL + lint/API restrictions
  \item[Automation] Yes/Partial
  \item[Sources] \texttt{NODE-SEC, GH-BUILTIN}
\end{description}
\end{rulebox}
\begin{rulebox}[title={\texttt{SCS-JS-INPUT-001} \hfill \prioritybadge{High Priority}{High}}]
\textbf{Domain.} Input Validation

\textbf{Requirement.} When copying or merging untrusted object properties, dangerous prototype-related keys MUST be rejected or the application MUST use data structures/APIs that do not mutate object prototypes.

\textbf{Rationale.} Prototype pollution can change application behavior globally and create downstream injection or authorization failures.

\textbf{Risk.} Prototype pollution

\begin{description}[leftmargin=2.65cm,labelwidth=2.3cm,style=nextline,itemsep=1pt,topsep=3pt]
  \item[CWE] CWE-1321, CWE-915
  \item[OWASP] Context-dependent / not forced
  \item[CERT] Not mapped
  \item[CodeQL] \textbf{Direct}; js/prototype-polluting-assignment
  \item[Enforcement] CodeQL + safe merge/object patterns + tests
  \item[Prevention] Schema/domain validation at trust boundaries before security-sensitive use.
  \item[Verification] CodeQL + safe merge/object patterns + tests
  \item[Automation] Yes
  \item[Sources] \texttt{NODE-SEC, GH-BUILTIN}
\end{description}
\end{rulebox}
\begin{rulebox}[title={\texttt{SCS-JS-INJ-001} \hfill \prioritybadge{Critical Control}{Critical}}]
\textbf{Domain.} Injection Prevention

\textbf{Requirement.} Node.js child processes MUST use executable/argument APIs without shell interpretation for untrusted arguments; shell strings MUST NOT contain untrusted data.

\textbf{Rationale.} Shell metacharacters can transform input into commands.

\textbf{Risk.} Command injection

\begin{description}[leftmargin=2.65cm,labelwidth=2.3cm,style=nextline,itemsep=1pt,topsep=3pt]
  \item[CWE] CWE-78
  \item[OWASP] Context-dependent / not forced
  \item[CERT] Not mapped
  \item[CodeQL] \textbf{Direct}; JavaScript command-injection queries
  \item[Enforcement] CodeQL + child\_process API policy
  \item[Prevention] Structured/parameterized APIs; keep untrusted data out of interpreter syntax.
  \item[Verification] CodeQL + child\_process API policy
  \item[Automation] Yes/Partial
  \item[Sources] \texttt{NODE-CHILD, GH-BUILTIN}
\end{description}
\end{rulebox}
\begin{rulebox}[title={\texttt{SCS-JS-WEB-001} \hfill \prioritybadge{Critical Control}{Critical}}]
\textbf{Domain.} Web Application Security

\textbf{Requirement.} Browser/server JavaScript MUST keep untrusted data out of unsafe HTML/DOM/script sinks and use framework/context-aware escaping APIs.

\textbf{Rationale.} DOM and server-template sinks can interpret data as markup or code.

\textbf{Risk.} XSS / DOM XSS

\begin{description}[leftmargin=2.65cm,labelwidth=2.3cm,style=nextline,itemsep=1pt,topsep=3pt]
  \item[CWE] CWE-79
  \item[OWASP] OWASP ASVS 5.0 --- output encoding
  \item[CERT] Not mapped
  \item[CodeQL] \textbf{Direct}; JavaScript XSS / unsafe DOM queries
  \item[Enforcement] CodeQL + framework-safe rendering + DAST
  \item[Prevention] Framework browser/session protections plus server-side validation and authorization.
  \item[Verification] CodeQL + framework-safe rendering + DAST
  \item[Automation] Yes/Partial
  \item[Sources] \texttt{OWASP-ASVS, GH-BUILTIN}
\end{description}
\end{rulebox}
\begin{rulebox}[title={\texttt{SCS-JS-RESOURCE-001} \hfill \prioritybadge{High Priority}{High}}]
\textbf{Domain.} Resource Management

\textbf{Requirement.} Node.js network services MUST bound request body sizes, timeouts, socket/connection use, regex complexity, and other attacker-controlled resource costs.

\textbf{Rationale.} The event loop and shared process resources can be monopolized by adversarial work.

\textbf{Risk.} Denial of service

\begin{description}[leftmargin=2.65cm,labelwidth=2.3cm,style=nextline,itemsep=1pt,topsep=3pt]
  \item[CWE] CWE-400, CWE-1333
  \item[OWASP] Context-dependent / not forced
  \item[CERT] Not mapped
  \item[CodeQL] \textbf{Partial}; JavaScript regex/resource queries
  \item[Enforcement] CodeQL + runtime limits + load/fuzz tests
  \item[Prevention] Explicit sizes, quotas, timeouts, backpressure, and complexity limits.
  \item[Verification] CodeQL + runtime limits + load/fuzz tests
  \item[Automation] Partial
  \item[Sources] \texttt{NODE-SEC, GH-BUILTIN}
\end{description}
\end{rulebox}
\begin{rulebox}[title={\texttt{SCS-JS-DEP-001} \hfill \prioritybadge{Critical Control}{Critical}}]
\textbf{Domain.} Dependency / Package Security

\textbf{Requirement.} Package installation in CI MUST use committed lockfiles and controlled registries; package lifecycle scripts and dependency updates SHOULD be reviewed according to project supply-chain policy.

\textbf{Rationale.} The npm ecosystem executes third-party code during build/install and has a large transitive dependency surface.

\textbf{Risk.} Supply-chain compromise

\begin{description}[leftmargin=2.65cm,labelwidth=2.3cm,style=nextline,itemsep=1pt,topsep=3pt]
  \item[CWE] CWE-829
  \item[OWASP] Context-dependent / not forced
  \item[CERT] Not mapped
  \item[CodeQL] \textbf{Requires another control}; ---
  \item[Enforcement] Dependabot + dependency review + lockfile-controlled CI
  \item[Prevention] Approved registries, reproducible dependency state, dependency review, and timely remediation.
  \item[Verification] Dependabot + dependency review + lockfile-controlled CI
  \item[Automation] Yes
  \item[Sources] \texttt{NODE-SEC, GH-DEP-REVIEW, GH-DEPENDABOT}
\end{description}
\end{rulebox}
\section{Secure / Insecure Implementation Example}
\subsubsection*{Representative implementation example}
\subsubsection*{Example --- avoid shell interpretation}

\textbf{Noncompliant}\par
\begin{lstlisting}[language={Java}]
exec(`git show ${userRef}`, callback);
\end{lstlisting}

\textbf{Compliant}\par
\begin{lstlisting}[language={Java}]
execFile("git", ["show", userRef], callback);
\end{lstlisting}

The secure form preserves executable/argument boundaries instead of asking a shell to parse an attacker-influenced command string. The application still needs business validation of which refs are permitted. Relevant weakness: CWE-78.\par
\begin{callout}{Navy}
\textbf{Example scope.} Examples illustrate selected high-risk patterns, not an exhaustive API allowlist. The trust boundary, sink, residual assumptions, CWE mapping, and CodeQL applicability must still be evaluated in the concrete application context.
\end{callout}
\section{Language-Specific CodeQL Query Crosswalk}\begin{landscape}\scriptsize
\setlength{\tabcolsep}{2.6pt}
\begin{longtable}{@{}P{1.28in}P{1.12in}P{1.52in}P{0.78in}P{0.52in}P{0.48in}P{1.32in}P{0.82in}@{}}
\toprule
\textbf{Rule ID(s)} & \textbf{Weakness} & \textbf{Query ID} & \textbf{CWE} & \textbf{Prec.} & \textbf{Sec.} & \textbf{Suite(s)} & \textbf{Coverage}\\
\midrule
\endfirsthead
\toprule
\textbf{Rule ID(s)} & \textbf{Weakness} & \textbf{Query ID} & \textbf{CWE} & \textbf{Prec.} & \textbf{Sec.} & \textbf{Suite(s)} & \textbf{Coverage}\\
\midrule
\endhead
SCS-ALL-INJ-001 & SQL injection & \href{https://codeql.github.com/codeql-query-help/javascript/js-sql-injection/}{\code{js/sql-injection}} & CWE-89, CWE-90, CWE-943 & high & 8.8 & code-scanning; sec.-extended; sec.-and-quality & Direct example\\
SCS-JS-INPUT-001 & Prototype-polluting assignment & \href{https://codeql.github.com/codeql-query-help/javascript/js-prototype-polluting-assignment/}{\code{js/prototype-polluting-assignment}} & CWE-78, 79, 94, 400, 471, 915 & high & 6.1 & code-scanning; sec.-extended; sec.-and-quality & Direct example\\
SCS-ALL-WEB-001 & Server-side request forgery & \href{https://codeql.github.com/codeql-query-help/javascript/js-request-forgery/}{\code{js/request-forgery}} & CWE-918 & high & 9.1 & code-scanning; sec.-extended; sec.-and-quality & Direct example\\
\bottomrule
\end{longtable}
\normalsize\end{landscape}
\section{Non-CodeQL Enforcement Matrix}
\begin{tabularx}{\textwidth}{@{}P{1.22in}Y P{1.55in} P{1.55in} P{1.5in}@{}}
\toprule
\textbf{Security area} & \textbf{Primary prevention} & \textbf{CodeQL} & \textbf{Companion automation / dynamic verification} & \textbf{Human review}\\
\midrule
Injection / data flow & Structured APIs, validation, contextual encoding & Primary static control where modeled & Unit/integration tests; DAST/IAST where web-facing & Trust-boundary review\\
Memory / concurrency & Safe abstractions, ownership, synchronization & Partial / language-dependent & runtime/framework linting; unit/integration tests; DAST/IAST for web apps & Unsafe/native/concurrency review\\
Secrets & External secret store; least privilege & Not primary & Secret scanning + push protection + rotation & Secret ownership / exception review\\
Dependencies & Approved sources; reproducible resolution & Not primary & Dependabot + dependency review + npm-compatible package ecosystem and lockfiles & Upgrade/exception governance\\
AuthN / AuthZ & Framework/IdP and server-side policies & Partial / framework-dependent & Authorization matrix + integration/DAST/IAST tests & Architecture/security review\\
Cryptography / TLS & Platform cryptography and trust validation & Partial API-misuse detection & Network/interoperability tests & Cryptographic design review\\
\bottomrule
\end{tabularx}

\section{Dependency and Package Security}
For JavaScript repositories, dependency security is governed through \textbf{npm-compatible package ecosystem and lockfiles} together with GitHub supply-chain controls. Projects MUST use approved registries/sources, SHOULD preserve deterministic or locked dependency state where the ecosystem supports it, MUST review vulnerable direct and transitive dependencies, and MUST address dependency-confusion or malicious-package risks in repository/build configuration.

\begin{callout}{Teal}
\textbf{GitHub enforcement.} Dependabot alerts/security updates and dependency review are the primary GitHub controls for vulnerable dependency introduction and remediation. CodeQL may identify vulnerable usage patterns in application code but is not the primary vulnerable-component inventory mechanism.
\end{callout}

Pull requests that change manifests, lockfiles, package sources, or build scripts SHOULD receive dependency review and SHOULD be merge-blocked when they introduce dependencies that violate the organization's documented vulnerability, provenance, or license policy.

\section{Secrets and Credential Handling}
Secrets MUST NOT be hard-coded into source, fixtures, build scripts, examples, or logs. Repository secret detection is primarily enforced through GitHub secret scanning and push protection, supplemented by approved secret stores, least-privilege credentials, and documented credential rotation/revocation procedures.

\begin{callout}{Critical}
\textbf{Exposure response.} Removing a secret from the current file or commit is not sufficient. A credential that reached repository history, an artifact, or a log MUST be evaluated for revocation/rotation because copies may persist outside the working tree.
\end{callout}

CodeQL is not the primary secret-scanning mechanism and a clean CodeQL result MUST NOT be interpreted as evidence that no credentials were committed.

\section{SSDLC Integration}
\small
\begin{tabularx}{\textwidth}{@{}P{1.12in}Y P{2.15in}@{}}
\toprule
\textbf{Stage} & \textbf{Engineering activity} & \textbf{Required evidence / gate}\\
\midrule
Design & Identify trust boundaries, sensitive data, authorization model, external dependencies, native/unsafe boundaries, and abuse cases. & Threat/architecture review where risk warrants.\\
Implementation & Apply applicable SCS-ALL and language overlay rules. & Code review; compiler/type/lint diagnostics.\\
Pull Request & Review code, dependency, workflow, and secret changes. & Dependency review; push protection; reviewer approval.\\
CodeQL & Run the configured language queries and required suites. & No unaccepted merge-blocking findings.\\
Build / Test & Run language-native security checks and negative-path tests. & CI status checks, race/sanitizer/fuzz results as applicable.\\
Security Testing & Exercise runtime behavior not reliably proven statically. & DAST/IAST/integration/property/fuzz testing as appropriate.\\
Merge Gate & Enforce repository/ruleset policy. & Required checks + approved exceptions only.\\
Deployment & Preserve secure configuration, credentials, TLS, and least privilege. & Deployment policy/configuration evidence.\\
Runtime Monitoring & Capture security-relevant events without secrets. & Central monitoring and alerting.\\
Remediation & Fix, validate, merge, and close findings with traceability. & Retest/rescan; time-bounded risk acceptance if unresolved.\\
\bottomrule
\end{tabularx}
\normalsize

\section{Developer Pre-Merge Checklist}
\begin{itemize}[leftmargin=*,itemsep=2pt]
\item Untrusted input is validated at each trust boundary.
\item Structured/parameterized APIs prevent interpreter injection.
\item Dynamic execution is avoided or strictly constrained.
\item Sensitive data and secrets are not logged or committed.
\item Filesystem paths, archive extraction, and temporary-file behavior are constrained.
\item TLS certificate/hostname verification remains enabled and outbound destinations are constrained where required.
\item Dependencies and lock/manifests were reviewed; vulnerable additions are resolved or formally accepted.
\item CodeQL findings and required companion security checks are resolved or covered by an approved, time-bounded exception.
\item Prototype pollution, DOM sinks, eval/Function, child processes, request limits, and npm lifecycle/dependency risks reviewed.
\item Required unit, integration, fuzz/property, race/sanitizer, DAST/IAST, or framework-security tests appropriate to the change have passed.
\end{itemize}

\section{Recommended Verification Controls}
The canonical catalog records the following language-specific verification profile:

CodeQL JavaScript/TypeScript + runtime/framework linting/tests + DAST/IAST for web apps + dependency review/Dependabot + secret scanning + lockfile/install-policy enforcement.

This profile is additive to the rule-level enforcement fields and to organization-specific required checks.
\begin{landscape}
\section{Official Reference Catalog}
All references below are first-party or authoritative sources used by this standalone standard. Access / verification date: 23 September 2026.

\scriptsize
\setlength{\tabcolsep}{3pt}
\begin{longtable}{@{}P{0.82in}P{1.23in}P{2.25in}P{1.7in}P{3.0in}@{}}
\toprule
\textbf{Key} & \textbf{Organization} & \textbf{Document} & \textbf{Applies to} & \textbf{URL}\\
\midrule
\endfirsthead
\toprule
\textbf{Key} & \textbf{Organization} & \textbf{Document} & \textbf{Applies to} & \textbf{URL}\\
\midrule
\endhead
APPLE-ATS & Apple & Preventing Insecure Network Connections & Swift / Apple network transport & \url{https://developer.apple.com/documentation/security/preventing-insecure-network-connections}\\
APPLE-CRYPTO & Apple & Apple CryptoKit & Swift / Apple cryptography & \url{https://developer.apple.com/documentation/cryptokit}\\
ASPNET-CSRF & Microsoft & Prevent Cross-Site Request Forgery (XSRF/CSRF) attacks in ASP.NET Core & C\# / ASP.NET Core & \url{https://learn.microsoft.com/en-us/aspnet/core/security/anti-request-forgery}\\
CERT-C & SEI / Carnegie Mellon University & SEI CERT C Coding Standard, 2016 Edition & C & \url{https://www.sei.cmu.edu/library/sei-cert-c-coding-standard-rules-for-developing-safe-reliable-and-secure-systems-2016-edition/}\\
CERT-CPP & SEI / Carnegie Mellon University & SEI CERT C++ Coding Standard, 2016 Edition & C++ & \url{https://www.sei.cmu.edu/library/sei-cert-c-coding-standard-rules-for-developing-safe-reliable-and-secure-systems-2016-edition-2/}\\
DOTNET-BF & Microsoft & Deserialization risks in use of BinaryFormatter and related types & C\# / .NET deserialization & \url{https://learn.microsoft.com/dotnet/standard/serialization/binaryformatter-security-guide}\\
DOTNET-CRYPTO & Microsoft & Cryptographic services & C\# / .NET cryptography & \url{https://learn.microsoft.com/en-us/dotnet/standard/security/cryptographic-services}\\
DOTNET-SEC & Microsoft & Secure coding guidelines for .NET & C\# / .NET & \url{https://learn.microsoft.com/en-us/dotnet/standard/security/secure-coding-guidelines}\\
GH-BUILTIN & GitHub & CodeQL built-in queries & Built-in CodeQL query inventories & \url{https://docs.github.com/en/code-security/reference/code-scanning/codeql/codeql-queries/built-in-queries}\\
GH-DEP-REVIEW & GitHub & About dependency review & Dependency review and pull-request gating & \url{https://docs.github.com/en/code-security/supply-chain-security/understanding-your-software-supply-chain/about-dependency-review}\\
GH-DEPENDABOT & GitHub & About Dependabot security updates & Vulnerable dependency remediation & \url{https://docs.github.com/en/code-security/dependabot/dependabot-security-updates/about-dependabot-security-updates}\\
GH-LANG & GitHub & CodeQL --- Supported languages and frameworks & All CodeQL languages; versions, compilers, frameworks, libraries & \url{https://codeql.github.com/docs/codeql-overview/supported-languages-and-frameworks/}\\
GH-PUSH & GitHub & About push protection & Blocking supported secrets before repository exposure & \url{https://docs.github.com/en/code-security/secret-scanning/introduction/about-push-protection}\\
GH-SECRET & GitHub & About secret scanning & Secret discovery & \url{https://docs.github.com/en/code-security/secret-scanning/introduction/about-secret-scanning}\\
GO-RAND & The Go Project & crypto/rand package & Go cryptographic randomness & \url{https://pkg.go.dev/crypto/rand}\\
GO-SEC & The Go Project & Security Best Practices for Go Developers & Go & \url{https://go.dev/doc/security/best-practices}\\
JAVA-26 & Oracle & JDK 26 Documentation --- Security & Java 26 security APIs and specifications & \url{https://docs.oracle.com/en/java/javase/26/}\\
JAVA-SCG & Oracle & Secure Coding Guidelines for Java SE, document version 11.0 (updated June 2025) & Java & \url{https://www.oracle.com/java/technologies/javase/seccodeguide.html}\\
MITRE-CWE & MITRE & Common Weakness Enumeration (CWE) & Weakness taxonomy and mappings & \url{https://cwe.mitre.org/}\\
NIST-SSDF & NIST & SP 800-218, Secure Software Development Framework (SSDF) Version 1.1 & Cross-language secure software development practices & \url{https://csrc.nist.gov/pubs/sp/800/218/final}\\
NODE-CHILD & OpenJS Foundation / Node.js & Node.js child\_process API & JavaScript process execution & \url{https://nodejs.org/api/child_process.html}\\
NODE-SEC & OpenJS Foundation / Node.js & Node.js Security Best Practices & Server-side JavaScript / Node.js & \url{https://nodejs.org/en/learn/getting-started/security-best-practices}\\
OWASP-ASVS & OWASP & Application Security Verification Standard (ASVS) 5.0 & Web/application security verification requirements & \url{https://owasp.org/www-project-application-security-verification-standard/}\\
PY-PICKLE & Python Software Foundation & Python 3.14 --- pickle & Python deserialization & \url{https://docs.python.org/3.14/library/pickle.html}\\
PY-TAR & Python Software Foundation & Python 3.14 --- tarfile & Python archive extraction & \url{https://docs.python.org/3.14/library/tarfile.html}\\
RAILS-SEC & Ruby on Rails & Ruby on Rails Security Guide & Ruby / Rails web security & \url{https://guides.rubyonrails.org/security.html}\\
RUBY-MARSHAL & Ruby Core & Ruby 3.3 Marshal --- Security considerations & Ruby deserialization & \url{https://ruby-doc.org/3.3/core/Marshal.html}\\
RUST-FFI & The Rust Project & The Rustonomicon --- FFI & Rust foreign-function interfaces & \url{https://doc.rust-lang.org/nomicon/ffi.html}\\
RUST-UNSAFE & The Rust Project & The Rustonomicon --- How Safe and Unsafe Interact & Rust unsafe boundaries & \url{https://doc.rust-lang.org/stable/nomicon/safe-unsafe-meaning.html}\\
SWIFT-INTEROP & The Swift Project & Safely Mixing Swift and C/C++ & Swift C/C++ interoperability & \url{https://www.swift.org/documentation/cxx-interop/safe-interop/}\\
SWIFT-MEM & The Swift Project & The Swift Programming Language --- The Basics: Memory Safety & Swift memory safety & \url{https://docs.swift.org/swift-book/documentation/the-swift-programming-language/thebasics/}\\
TS-STRICT & Microsoft / TypeScript & TSConfig --- strict & TypeScript strict type checking & \url{https://www.typescriptlang.org/tsconfig/strict}\\
\bottomrule
\end{longtable}
\normalsize
\end{landscape}

\appendix
\section{Build and Maintenance}
Compile this standalone standard with:
\begin{lstlisting}
latexmk -lualatex -interaction=nonstopmode -halt-on-error secure-coding-standard-javascript.tex
\end{lstlisting}
Revalidate the CodeQL support profile, query metadata, language/platform release information, and authoritative guidance whenever the CodeQL bundle or language/toolchain baseline changes.

\section{Rule Identifier Semantics}
Universal controls retain `SCS-ALL-<DOMAIN>-<NNN>`. Language overlays retain their original `SCS-<LANG>-<DOMAIN>-<NNN>` identifiers. Universal controls are not renamed when copied into this standalone standard.
\end{document}
